% 書末：參考資料與延伸閱讀（不編章號）
\chapter*{參考資料與延伸閱讀}
\addcontentsline{toc}{chapter}{參考資料與延伸閱讀}
\markboth{參考資料與延伸閱讀}{參考資料與延伸閱讀}
\label{ch:references}

各章最後的「延伸閱讀」提到的書，完整的書目資料都列在這裡，每本附一兩句它適合
做什麼，以及和本書哪幾章搭配。 括號裡的章節號碼是那本書的章節，以這裡列出的
版本為準；換了版本，章節號碼可能不同。

本書的講解、例子和程式都是自己寫的，參考這些書的地方是題型的分類、講解的先後
和經典的結論。 想看另一種講法、更嚴謹的證明，或想多做一些題目時，可以從下面
挑一本。

\section*{篇幅短、適合快速複習}

\begin{itemize}
\item Antti Laaksonen，\textit{Competitive Programmer's Handbook}，2018 年 7 月
3 日草稿版，作者公開的 PDF。

每個主題只用幾頁、一兩個小例子講完，適合讀完本書某一章後，花十幾分鐘從另一個
角度再看一次。 和本書搭配的地方：第~1、2~章（第 7 章的背包、最長遞增子序列、
編輯距離），第~6~章（第 16 章前半），第~7~章（12.3 節，第 15、16 章），第~8~章
（第 13 章），第~9~章（第 15 章），第~10~章（第 14 章）。

\item Antti Laaksonen，\textit{Guide to Competitive Programming: Learning and
Improving Algorithms Through Contests}，第 3 版，Springer，2024。

和上一本出自同一位作者，章節安排相近，是正式出版的書，內容較新。 本書第~4~%
章引用了它 6.2.4、6.2.5 節：從排列改成子集合的動態規劃，以及一列一列數鋪磚方法
的寫法。
\end{itemize}

\section*{題型分類細、題目多}

\begin{itemize}
\item Steven Halim、Felix Halim、Suhendry Effendy，\textit{Competitive
Programming 4}，Book 1 與 Book 2，2020（作者自行出版）。

把題目依技巧分得很細，每一節後面附一長串分好類的練習題，適合打算參加程式競
賽、想大量練習的讀者。 和本書搭配的地方：第~1~章（Book 1 的 3.5.2 節），第~2~章（Book 1
的 3.5.2 節，Book 2 的 6.3 節），第~4~章（Book 1 的 3.5.2 節，Book 2 第 8 章），第~5~章
（Book 2 的 6.3.3 節），第~6~章（Book 1 的 4.2.6 節），第~7~章（Book 1 的
2.4.2、4.2.7、4.2.8 節），第~8~章（Book 1 的 4.4、4.5 節），第~9~章（Book 1
的 4.3 節）。

\item 劉汝佳，《算法競賽入門經典》，第 2 版，清華大學出版社，2014。

簡體中文。 例題多，程式寫得精簡，適合想用中文多看幾種題型的讀者。 和本書搭配
的地方：第~1~章（9.3.2 節），第~3~章（9.4.1 節），第~5~章（10.4 節例題
10-22），第~6~章（6.4.3、9.2 節），第~8~章（11.3 節），第~10~章（9.4.2 節）。

\item 吳邦一，《AP325：從 APCS 實作題檢測三級到五級》，v1.5，2024 年 9 月，作者
公開的講義。

以 APCS 實作題的程度為目標，每個主題附很多例題和習題，還有測試資料，適合要準
備 APCS 的讀者多練。 它的第 6 章（動態規劃）、第 7 章（基本圖論演算法）、第 8
章（樹上演算法）和本書的範圍重疊最多；各章引用的是第~3~章（6.2.4 節）和第~10~%
章（8.3 節）。
\end{itemize}

\section*{嚴謹的證明與演算法設計}

\begin{itemize}
\item Thomas H. Cormen、Charles E. Leiserson、Ronald L. Rivest、Clifford Stein，
\textit{Introduction to Algorithms}，第 4 版，The MIT Press，2022。

大學演算法課常用的教科書，每個演算法都有完整的正確性證明和分析。 想知道本書
「為什麼一定對」框裡省略的細節，或要修演算法課的讀者，可以對照著讀。 和本書搭
配的地方：第~2~章（14.4 節），第~3~章（14.2 節），第~6~章（20.4、22.2 節），
第~7~章（第 19 章），第~8~章（第 22 章、23.2 節），第~9~章（第 21 章）。

\item Steven S. Skiena，\textit{The Algorithm Design Manual}，第 3 版，Springer，
2020。

重點放在怎麼想出演算法、怎麼把實際的問題對應到已知的問題，後半是一份常見問題
的目錄。 適合想練習建模的大學生和自學的人。 和本書搭配的地方：第~2~章（10.2
節），第~4~章（10.9 節）。
\end{itemize}

\section*{本系列的其他書}

\begin{itemize}
\item 《C++ 演算法解題 1：APCS 實作題核心》：本書的前一本，各章提到的「第 1 本」。
排序、二分搜、前綴和、堆疊與佇列、貪心、一維與格子的動態規劃、BFS 與 DFS、樹的
基本做法都在那裡。
\item 《C++ 程式設計：給學過 C 的人》：各章提到的「銜接書」，本書用到的 C++ 語法
和 STL 都在那裡。
\end{itemize}
